\section{入力：KN のデータを超える refined 不変量}\label{sec:inputs}

\subsection{トーラス結び目}
トーラス結び目 $T_{m,n}$ の $\overline{\rCS}_\lambda$ は，Gorsky--Neguţ~\cite{Gorsky13b} による楕円 Hall 代数の生成元の作用で計算する．
楕円 Hall 代数と DAHA の対応は Schiffmann--Vasserot~\cite{Schiffmann09} により，ケーブルの場合の対応は Morton--Samuelson~\cite{Morton14} による．
% [I1]
規約を固定するため，KN~\cite{Kameyama17} の補助ファイルにあるトーラス結び目の値 27 件（$T_{2,3}$, $T_{2,5}$, $T_{2,7}$, $T_{2,9}$ の $|\lambda|\le3$ と $T_{3,4}$ の $\yng(1),\yng(2),\yng(1,1)$）と，
KN の付録 C の BPS の表 13 個を，単項式の正規化（\S\ref{sec:normalization}）の下で全係数で再現した．
これは既知の結果の再現であり，以下の計算の規約の確認として用いる．

\subsection{代数的な iterated torus 結び目}
companion $K$ の $(m,n)$-ケーブルを $C_{m,n}K$ と書く．$m$ は巻き数で，$n$ は companion の Seifert framing で測る．
Cherednik--Danilenko~\cite{Cherednik14,Cherednik15} の DAHA による定義では，$C_{m_2,n_2}T_{m_1,n_1}$ の Newton 対は $r=(m_1,m_2)$，$s=(n_1,n_2-m_1n_1m_2)$ であり，
全ての $r_i,s_i>0$（$n_2>m_1n_1m_2$）のとき代数的である．
% [I2]
我々はこの定義を companion の段で線形に分解し，楕円 Hall 代数の演算子の合成として計算した．
すなわち一番外側の段 $C_{m,n}$ について $\Phi_\lambda=P_\lambda[p_k\mapsto P_{k(m,-s')}]\cdot1=\sum_\mu c_{\lambda\mu}(s')P_\mu$ と展開し（$s'$ は Newton 対 $s$ の最後の成分，$p_k$ は冪和），
companion の値 $\rCS_K(\mu)$ との一次結合をとる．
素朴な Rosso--Jones の $t$ 変形が誤りであることは知られており~\cite{Danilenko14,Morozov18}，ここで使う形は定義から直接従うものである．
\begin{compresult}\label{cr:daha}
この計算法は，Cherednik--Danilenko の明示例 10 件（$K1$ の $\yng(1)$ と $\yng(2)$，$K2,K3,K4,K5a,K5b,K5c,K7,K8$ の $\yng(1)$．非代数的なケーブルの名前は付録~\ref{app:nonalg}）と，単項式倍を除いて全係数で一致する．
また，$t=q$ の極限はブレイドから計算した色付き HOMFLY-PT 多項式と一致する．
\end{compresult}

本論文で扱うケーブルを表~\ref{tab:cables} にまとめる．
\begin{table}[h]
\centering
\caption{代数的な iterated torus 結び目．$C_{m,n}K$ は $K$ の $(m,n)$-ケーブル．$b$ はブレイドの本数，$\delta$ は種数型の量（\S\ref{sec:normalization}）．}\label{tab:cables}
\begin{tabular}{llcc}
\toprule
名前 & 結び目 & $b$ & $\delta$ \\
\midrule
$K1$ & $T_{2,3}$ の $(2,13)$-ケーブル & 4 & 8 \\
$K3$ & $T_{2,3}$ の $(2,15)$-ケーブル & 4 & 9 \\
$K6a$ & $T_{2,5}$ の $(2,21)$-ケーブル & 4 & 14 \\
$K9$ & $T_{2,7}$ の $(2,31)$-ケーブル & 4 & 21 \\
$K8$ & $T_{3,4}$ の $(2,25)$-ケーブル & 6 & 18 \\
$K7$ & $T_{2,3}$ の $(3,19)$-ケーブル & 6 & 21 \\
$K10$ & $K1$ の $(2,53)$-ケーブル & 8 & 42 \\
\bottomrule
\end{tabular}
\end{table}

\subsection{単項式の正規化}\label{sec:normalization}
DAHA と楕円 Hall 代数は値を単項式倍を除いてしか決めないが，KN の BPS 抽出は単項式倍で結果が変わる．
そこで次の規則で正規化する．
\begin{itemize}
\item[(N1)] $t=q$ の極限が，framing 0 の reduced 色付き HOMFLY-PT 多項式（DGR の規約）に等しい．
\item[(N2)] $\yng(1)$ の太字の変数での最低の $\bt$ 次数が 0 である．
\item[(N2$'$)] 表現 $\lambda$ の単項式は $M_\lambda=M_\square^{|\lambda|}\,T_\lambda^{F}$ である．$F$ は一番外側のケーブル $C_{m,n}$ では $mn$，トーラス結び目 $T_{m,n}$ では $mn$ である．
\end{itemize}
トーラス結び目では，(N1)(N2) の下で KN の値 27 件が全係数で再現される．
(N2) が決めるのは，箱あたりの $(q/t)^{c}$ の自由度だけである．補題~\ref{lem:adams} により，これは正値性に影響しない．
色付きの場合の単項式 $\pm a^\alpha q^\beta t^\gamma$ のうち，$t=q$ で決まるのは符号，$\alpha$，$\beta+\gamma$ である．
色付きの場合の (N1) は検査として確かめた．
$t=q$ で決まらないのは，表現ごとの $(q/t)^{c_\lambda}$ の自由度だけである．

% [I4]
\begin{hypothesis}\label{hyp:N2}
ケーブルについて，$t=q$ で決まらない表現ごとの単項式 $(q/t)^{c_\lambda}$ は (N2$'$) の形をとる．
すなわち，$|\lambda|$ に比例する部分（正値性に無関係）と，refined の framing の固有値 $T_\lambda^F$ の分だけである．
\end{hypothesis}
命題~\ref{prop:N2-framing} により，仮説~\ref{hyp:N2} は「表現によって変わる部分は Seifert framing の $F$ 単位の framing である」と言い換えられる．
根拠は次の三つである．
(i) 後述の refined 衛星公式（仮説~\ref{hyp:sat}）から従う．
(ii) 非代数的なケーブル $K5c$ では，$\yng(2)$ を $(q/t)^u$ 倍する 1 パラメータ族のうち，Laurent 性と転置対称性を満たすのは (N2$'$) の $u$ だけである．
(iii) 9 個のケーブルの全ての表現で，(N2$'$) の形が計算した値と一致する．

% [I5]
\begin{compresult}\label{cr:M}
(N1)(N2) で決めた $\yng(1)$ の単項式は，計算した全 22 結び目（トーラス結び目 9 個，ケーブル 13 個）で
\[
M_\square=(-q/t)^{b-1}\bigl(a(q/t)^{1/2}\bigr)^{\delta}
\]
に一致する．ここで $b$ はブレイドの本数，$\delta$ はケーブル $C_{m,n}$ の段ごとに $\delta\mapsto m\delta+(m-1)(n-1)/2$ で決まる（$T_{m,n}$ では $\delta=(m-1)(n-1)/2$）．
\end{compresult}

% [I6]
\begin{hypothesis}[refined 衛星公式]\label{hyp:sat}
$\rCS_{C_{m,n}K}(\lambda)=\sum_\mu C_{\lambda\mu}\,\rCS_K(\mu)$ が成り立つ．
係数は，トーラス結び目 $T_{m,n}$（自明な結び目の $(m,n)$-ケーブル）の展開係数を (N2$'$) で正規化した $C_{\lambda\mu}=M_{T_{m,n}}^{|\lambda|}T_\lambda^{mn}c_{\lambda\mu}(n)$ である．
\end{hypothesis}
$t=q$ では仮説~\ref{hyp:sat} は古典的な衛星公式に帰着し，(N1) と両立する．

導出には，計算で確かめた次の二つの式を用いる．
\begin{compresult}\label{cr:tau}
framing 演算子 $\tau:P_\mu\mapsto T_\mu P_\mu$ について，楕円 Hall 代数の傾き $(m,n)$ の元と整数 $j$ に対し
$\tau^jP_\lambda[p_k\mapsto P_{k(m,n)}]\cdot1=P_\lambda[p_k\mapsto P_{k(m,n+jm)}]\cdot1$ が成り立つ（5 つの傾きと $j=1,-1,2$ で，有理数の点で厳密に一致）．
\end{compresult}
\begin{compresult}\label{cr:mellit}
Mellit~\cite[Thm.~3.1]{Mellit17} の Dyck 路の表示によるトーラス結び目の HOMFLY-PT ホモロジーの Poincaré 多項式 $\mathcal{P}_{m,n}$ は，
\[
\mathcal{P}_{m,n}=-\bigl(a(qt)^{-1/2}\bigr)^{\delta}\frac{1}{1-t}\,\frac{(Q_{m,n}\cdot1)[a/q-1]}{a/q-1}
\]
に一致する（$Q_{m,n}$ は Bergeron--Garsia--Leven--Xin~\cite{Bergeron14} の演算子．$(m,n)=(2,3),(3,2),(2,5),(3,4),(4,3),(3,5),(5,3)$ で厳密に一致）．
\end{compresult}

\begin{proposition}\label{thm:M}
仮説~\ref{hyp:sat} と計算結果~\ref{cr:tau}，\ref{cr:mellit} を仮定すると，計算結果~\ref{cr:M} の閉じた式と (N2$'$) の形が成り立つ．
\end{proposition}
\begin{proof}[証明の概略]
ケーブルの部分では，計算結果~\ref{cr:tau}（framing 演算子が傾きを $(m,n)\mapsto(m,n+jm)$ に送ること）を使う．
これと仮説~\ref{hyp:sat} から，$M_{C_{m,n}K}=M_K^m\,M_{T_{m,n}}$ と $F=mn$ が従う．
トーラス結び目の部分は四つの事実からなる．
(I) Gorsky--Neguţ の正規化前の値の最低 $a$ 係数は，有理 $q,t$-Catalan 数で与えられる．これは Gorsky--Neguţ の予想であり，Mellit~\cite{Mellit16} の定理の系として成り立つ．
(II) 最低 $a$ 次数は，Franks--Williams の定理による．
(III) $t=q$ での符号と $q$ の冪は，$q\leftrightarrow q^{-1}$ 対称性から決まる．
(IV) $(q/t)$ の指数は，Mellit の Dyck 路の表示（計算結果~\ref{cr:mellit}）から決まる．
詳細は付録~\ref{app:normalization} に述べる．
\end{proof}

\begin{remark}[Mellit の系の正規化]
Mellit~\cite[Cor.~3.4]{Mellit17} は $\mathcal{P}_{m,n}=(-1)^n(a(qt)^{-1/2})^\delta\frac{1}{1-t}\frac{P_{m,n}(1)[a-1]}{a-1}$ と書き，$P_{m,n}$ を「対応する楕円 Hall 代数の生成元」とする．
同論文の Thm.~3.3 から，この式は $P_{m,n}(1)=-d_-B_{c_{m,n}}d_+(1)$（Carlsson--Mellit の $d_\pm$ の表現の元）と読めば成り立つ．
一方，標準の正規化（Bergeron--Garsia--Leven--Xin の $Q_{m,n}$，Gorsky--Neguţ の $\widetilde P$）の元と読むと，代入 $[a-1]$ では成り立たない．
標準の正規化の元と変数の取り方を変えて調べると，成り立つのは計算結果~\ref{cr:mellit} の形だけである．
命題~\ref{thm:M} の (IV) が決めるのは $M_\square$ の $(q/t)$ の指数だけなので，この点は正値性の結果に影響しない（補題~\ref{lem:adams}）．
\end{remark}

\begin{remark}[Cherednik--Danilenko の tilde 正規化との関係]\label{rem:tilde}
Cherednik--Danilenko~\cite{Cherednik14} は DAHA superpolynomial を，最低の $a$ の冪と $q_{st}^\bullet t_{st}^\bullet$ で割って，$a_{st}=0$ の部分の定数項が 1 になるように正規化する（tilde 正規化）．
標準の変数 $(q_{st},t_{st},a_{st})$ は太字の変数と同じである．
係数が非負の計算値 19 個（トーラス結び目 9 個，ケーブル 10 個）で，本論文の $\yng(1)$ の太字の $P$ は
\[
P=\ba^{2\delta}\bq^{-2\delta}\times(\text{tilde 正規化})
\]
である（計算）．Cherednik--Danilenko の明示例との比も，正のもの（$K1,K3,K5b,K7,K8$）で全て $\ba^{2\delta}\bq^{-2\delta}$ である．
これらの $P$ は全て $\sigma$ で不変で，$\bq$ の指数は偶数である．
\end{remark}

\subsection{ホモロジーを入力とする結び目}
11 交点以下の全素結び目 801 個について，reduced HOMFLY-PT ホモロジーの Poincaré 多項式（Nakagane--Sano~\cite{Nakagane21} のデータ）を $\yng(1)$ の入力とする．
DGR の規約への変換は，$\sigma$ 不変性，$\bt=-1$ での HOMFLY-PT 多項式，トーラス結び目の $\rCS_{\yng(1)}$ との一致で確かめた．

\subsection{2 箱の予想式}
双曲結び目の refined な $\yng(2)$, $\yng(1,1)$ は定義どおりには得られない．そこで文献の予想式を入力とする．
\begin{itemize}
\item ダブルツイスト結び目 $K_{m,n}$ とツイスト結び目 $K_p=K_{p,1}$：Kameyama--Nawata--Tao--Zhang（KNTZ）~\cite{Kameyama19} の，長方形のヤング図で色付けた HOMFLY-PT ホモロジーの Poincaré 多項式の式（refined 版，予想）を用いる．
KNTZ の $K_{m,n}$ は，GNSSS~\cite{Gukov15} の $(2m-1,1,2n-1)$-pretzel 結び目 $K_{s,t}$（$s=m$，$t=n$）と同じ結び目である．
KNTZ の式は GNSSS の四重次数の予想式と一致する（$\yng(1)$，$\yng(2)$，$\yng(1,1)$ について．下の計算結果~\ref{cr:kntz}）．
$t=q$ では KNTZ の式は Fang--Zhou~\cite{Fang26} により証明されている．
KN の $t_c$ 次数への変換は，Gorsky--Gukov--Stošić~\cite{Gorsky13} の次数の関係 $\mathbf{Q}\mapsto\bq^R$，$\mathbf{t}_r\mapsto\bq^{-1}$，$\mathbf{t}_c\mapsto\bt\bq$（$R$ は長方形の行数）による．
変換は，三葉結び目 $K_{1,1}=T_{2,3}$ の $\yng(1)$，$\yng(2)$，$\yng(1,1)$，$\yng(2,2)$ が楕円 Hall 代数の値と全係数一致することで較正した（Nawata ら~\cite{Nawata12} も参照）．
\item thick な結び目 $\mathbf{10_{139}}$，$\mathbf{10_{152}}$：Stošić~\cite{Stosic24} の quiver 型の式と，Kucharski--Reineke--Stošić--Sułkowski（KRSS）~\cite{Kucharski17b} の予想式で $t_c$ 次数を作る．
この方法は，$\mathbf{8_{19}}=T_{3,4}$ で楕円 Hall 代数の値と $\yng(2)$ から $\yng(1,1,1)$ まで全項一致する．
\end{itemize}
\begin{compresult}\label{cr:kntz}
KNTZ の refined の式と GNSSS の四重次数の式は，$K_{m,n}$ の $\yng(1)$，$\yng(2)$ と，ミラー対称性 $\mathbf{t}_r\leftrightarrow\mathbf{t}_c$ を介した $\yng(1,1)$ で，
符号の四つの象限それぞれの $|m|,|n|\in\{1,2,3\}$ で全係数一致する（144 件）．
両方の式は $m,n$ のそれぞれについて（象限ごとに）根が 3 個以下の指数和なので，全ての $(m,n)$ で一致する．
また KNTZ の式の実装は，KNTZ~\cite{Kameyama19} の補助ファイルの値と全係数一致する．
\end{compresult}
これらを入力とする主張は「予想式を入力とした場合の正値性」であり，定義どおりの入力による 3 箱の結果（\S\ref{sec:three}）とは確からしさが異なる．

\subsection{厳密な計算}
$\overline{\rCS}_\lambda$ と $\hf_\rho$ は，有理数の点ではなく，31 ビットの素数を法とする評価と補間で求める．
係数が大きいときは 2〜3 個の素数と中国剰余定理を用いる．
得られた多項式は，補間に使っていない素数の点で検証する．
最大の例 $K10$ の 3 箱は，spin 表示で約 183 万項，係数の最大は約 $1.1\times10^{19}$ である．
計算時間とメモリは付録~\ref{app:computation} にまとめる．
